\documentclass[10pt,
%superscriptaddress,
%groupedaddress,
%unsortedaddress,
%runinaddress,
%frontmatterverbose, 
%preprint,
%preprintnumbers,
%nofootinbib,
%nobibnotes,
%bibnotes,
 amsmath,amssymb,
 aps,prb, twocolumn
]{revtex4-2}

% Graphics and math
\usepackage{empheq}
\usepackage{graphicx}
\usepackage{bm}
\usepackage{dcolumn}
\usepackage{physics}
\usepackage{soul}
\usepackage{dsfont}
\usepackage{xcolor} % for more colors
\usepackage[percent]{overpic} % for captions
% Subfigures (compatible with REVTeX)
\usepackage[caption=false]{subfig}  
%\captionsetup[subfloat]{position=top,justification=raggedright,singlelinecheck=false}
% Hyperlinks (must come last)
\usepackage[colorlinks=true,allcolors=blue]{hyperref}


\newcommand{\id}{\mathds{1}}
\newcommand\sj[1]{ {\color{cyan} #1} } 
\newcommand\jm[1]{ {\color{magenta} #1} } 
\definecolor{mediumgreen}{HTML}{228B22} % forest green-ish
\newcommand\oz[1]{ {\color{mediumgreen} #1} } 

\begin{document}
%=====================================================================
%  APPENDIX B  (replacement)
%  Von Neumann's measurement scheme on a lattice
%=====================================================================
\section{TEST for Von Neumann's Measurement Scheme on a Lattice}
\label{app:vn}

Von Neumann's scheme idealizes measurement as an instantaneous unitary
kick~\cite{VonNeumann}. In this appendix we construct the lattice
counterpart appropriate to Eq.~(1), identify the operator that plays the
role of the continuum momentum, and extract the backaction phase
$\Delta\phi$ and the coherence factor $R$. Our central result is that the
detector coupling must be a function of the crystal momentum, and that the
physically correct function is fixed by the on-shell scattering phase of the
bond defect rather than chosen by fiat.

%---------------------------------------------------------------------
\subsection{Two conditions for a measurement kick}
\label{app:vn:conditions}

A kick is a single unitary $\hat U(\delta\tau)=e^{-i\hat h_I\delta\tau}$
applied at one instant. Reducing a genuine interaction, spread over the
window $\delta\tau$, to such a form requires that the time-ordered
exponential collapse. Hence the coupling must be a constant of the free
motion, i.e.\ $[\hat h_I,\hat H_0]=0$ with $\hat H_0=\hat H_D+\hat H_S$.
That single requirement splits into two physically distinct conditions:
\begin{align}
[\hat h_I,\hat H_S] &= 0 &&\text{(qubit: non-demolition)}, \label{eq:condS}\\
[\hat h_I,\hat H_D] &= 0 &&\text{(detector: kick validity)}. \label{eq:condD}
\end{align}
We break Eq.~\eqref{eq:condS} deliberately. That breaking is the subject of
this work. Equation~\eqref{eq:condD} is a separate matter. It carries no
information about the qubit and must hold for the kick description to exist
at all.

The interaction of Eq.~(2c) fails Eq.~\eqref{eq:condD}. The bond operator
$\hat c^{\dagger}_{b+1}\hat c_b+\mathrm{h.c.}$ does not commute with the
uniform hopping $\hat H_D$, so in the interaction picture it acquires an
explicit time dependence over the window. The failure is not technical, it
is physical: the electron does not reside at the bond, it flows through it.
No instant exists at which the interaction can be said to act.

%---------------------------------------------------------------------
\subsection{The crystal momentum and its functions}
\label{app:vn:K}

We replace the finite chain by an infinite one, for which the crystal
momentum $\hat k$ is well defined and
\begin{equation}
\hat H_D=-2t\cos(a\hat k).
\label{eq:HDk}
\end{equation}
Equation~\eqref{eq:HDk} settles the classification. Since $\hat H_D$ is a
function of $\hat k$ alone, every operator satisfying Eq.~\eqref{eq:condD}
is itself a function of $\hat k$, up to the residual $k\to-k$ degeneracy of
the cosine band. The most general detector coupling compatible with a
measurement kick is therefore
\begin{equation}
\hat h_I\,\delta\tau=\hat\Phi(\hat k)\otimes\ket{1}\bra{1},
\label{eq:hIgeneral}
\end{equation}
with $\hat\Phi$ a real function on the Brillouin zone. We stress that
Eq.~\eqref{eq:hIgeneral} is a classification and not an ansatz. Von
Neumann's continuum coupling $\hat h_I\propto\hat p$ is the single element
$\Phi(k)\propto k$ of this family.

Two properties of $\hat\Phi(\hat k)$ carry all of the physics. First, its
value at the carrier momentum sets the conditional phase. Second, its
derivative sets the conditional displacement of the pointer,
\begin{equation}
\lambda(k)=\partial_k\Phi(k),
\label{eq:lambda}
\end{equation}
which is the Wigner delay length. For a wavepacket
$\ket{\psi}=\int\!\mathrm{d}k\,\tilde\psi(k)e^{-ikx_c}\ket{k}$ the
exponential of a \emph{linear} $\Phi$ acts as an exact rigid translation,
\begin{equation}
e^{i\lambda\hat k}\ket{\psi(x_c)}=\ket{\psi(x_c-\lambda)},
\label{eq:translate}
\end{equation}
because each Fourier component is multiplied by $e^{i\lambda k}$ and the
weights $\tilde\psi(k)$ are untouched. Equation~\eqref{eq:translate} is
exact on the lattice, not a continuum limit, provided the packet does not
wrap the Brillouin zone: $\lambda\ll\pi/k_0$, which weak coupling
guarantees. We note that $\hat k$ itself has no such simple action on site
states, viz.\ $\hat k\ket{j}=(2ia)^{-1}(\ket{j+1}-\ket{j-1})$. Only the
exponential translates.

Moving on, a curved $\Phi$ both shifts and deforms the packet. That
distinction is the whole content of the present model. Von Neumann's
pointer translates rigidly. Ours translates and distorts, and the
distortion is governed by $\Phi''$.

%---------------------------------------------------------------------
\subsection{Physical identification of \texorpdfstring{$\Phi(k)$}{Phi(k)}}
\label{app:vn:Phi}

The function $\Phi$ is not free. Before and after the encounter the
wavepacket is asymptotically free, so all that it retains of the
interaction is the on-shell scattering data. For a single channel and a
single defect that data is one phase per momentum. Hence
$\Phi(k)=\arg\tau(k)$, with $\tau(k)$ the transmission amplitude of the
modified bond. Replacing the local defect by $\hat\Phi(\hat k)$ is
therefore a rewriting of the same physics in the only representation in
which the kick picture is consistent, and not the introduction of a
different model.

A transparent derivation follows from first-order perturbation theory. The
qubit state $\ket{1}$ shifts the bond amplitude by $\delta t=g/2$, and for a
Bloch state on a ring of $L$ sites the induced energy shift is
\begin{equation}
\delta E_k=\bra{k}\delta t\big(\ket{b\!+\!1}\bra{b}+\mathrm{h.c.}\big)\ket{k}
=\frac{2\delta t}{L}\cos(ka).
\label{eq:dE}
\end{equation}
The electron spends a time $L/v_g$ in the region, with $v_g=2ta\sin(ka)$,
so the accumulated phase is the product of energy shift and dwell time:
\begin{equation}
\Phi(k)=-\,\delta E_k\,\frac{L}{v_g}
=-\frac{g}{2t}\cot(ka).
\label{eq:Phi}
\end{equation}
The length $L$ cancels, as it must. Equation~\eqref{eq:Phi} is the lattice
analogue of the semiclassical $\int\!\delta E\,\mathrm{d}t$, and it makes
the momentum dependence intuitive. The denominator $\sin(ka)\propto v_g$
gives the familiar threshold enhancement: a slow electron dwells longer and
accumulates more phase. The numerator $\cos(ka)$ is the bond matrix element,
which vanishes at $k=\pi/2a$ because there the amplitudes on adjacent sites
are in quadrature, so a real change of the hopping produces no first-order
energy shift. Hence a band-centre electron is blind to a weak bond defect.

The exact amplitude confirms Eq.~\eqref{eq:Phi} and delimits its validity.
Solving the scattering problem for a chain with hopping $t$ everywhere
except $t'=t-g/2$ on the bond $(b,b+1)$, and writing $\eta=t'/t$, we obtain
\begin{equation}
\tau(k)=\frac{-2i\eta\sin (ka)}{e^{-ika}-\eta^{2}e^{ika}},
\qquad
r(k)=\frac{(\eta^{2}-1)e^{ika}}{e^{-ika}-\eta^{2}e^{ika}}.
\label{eq:tauexact}
\end{equation}
Expanding Eq.~\eqref{eq:tauexact} for $g/2t\ll1$ recovers
Eq.~\eqref{eq:Phi} together with the leading amplitude loss,
\begin{equation}
\arg\tau=-\frac{g}{2t}\cot(ka)\Big[1+\frac{g}{4t}\Big],
\qquad
|\tau|=1-\frac{1}{2}\Big(\frac{g}{2t}\Big)^{2}\frac{1}{\sin^{2}(ka)}.
\label{eq:expand}
\end{equation}
The phase is first order in $g$ and the amplitude loss is second order, in
agreement with the perturbative hierarchy of Appendix~C.

%---------------------------------------------------------------------
\subsection{The kick and the reduced density matrix}
\label{app:vn:kick}

We now apply the kick. The composite state at the moment the packet reaches
the bond is $\ket{\psi(x_b)}\otimes\ket{q}$ with
$\ket{q}=\cos(\theta_b/2)\ket{0}+\sin(\theta_b/2)e^{i\phi_0}\ket{1}$.
Conditioned on $\ket{1}$ the packet is scattered, so
\begin{equation}
\ket{\Psi'}=\cos\tfrac{\theta_b}{2}\ket{\psi}\otimes\ket{0}
+\sin\tfrac{\theta_b}{2}e^{i\phi_0}
\big[\tau(\hat k)\ket{\psi_\rightarrow}+r(\hat k)\ket{\psi_\leftarrow}\big]\otimes\ket{1},
\label{eq:Psiprime}
\end{equation}
where $\ket{\psi_\rightarrow}$ and $\ket{\psi_\leftarrow}$ denote the
forward and backward moving components. Tracing out the QPC gives a
reduced qubit state whose populations are unchanged, since the scattering is
elastic, and whose coherence is multiplied by the overlap
\begin{equation}
\big\langle\tau\big\rangle
\equiv\bra{\psi}\tau(\hat k)\ket{\psi}
=\int\!\mathrm{d}k\,|\tilde\psi(k)|^{2}\,\tau(k).
\label{eq:overlap}
\end{equation}
Only the forward component contributes: the reflected packet is
orthogonal to the undisturbed one, so the which-path information it carries
is lost to the trace. Equation~\eqref{eq:overlap} is exact and replaces the
Gaussian-displacement estimate of the previous treatment.

The momentum weight $|\tilde\psi(k)|^{2}$ is a normalized Gaussian of
variance $\sigma_k^{2}=1/(2\Delta_0^{2})$. We emphasize that $\sigma_k$ is
conserved under free evolution, whereas the real-space width disperses.
Hence the initial width $\Delta_0$, and not the width $\Delta(k_0)$ at the
time of interaction, controls the overlap. Evaluating
Eq.~\eqref{eq:overlap} by a second-order cumulant expansion of
$\log\tau=\log|\tau|+i\Phi$ yields
\begin{equation}
\big\langle\tau\big\rangle
=R\,e^{-i\Delta\phi},
\label{eq:Rphi}
\end{equation}
with
\begin{empheq}[box=\fbox]{align}
\Delta\phi(k_0)&=\frac{g}{2t}\cot(ak_0)
+\frac{g}{4t}\frac{\cos(ak_0)}{\Delta_0^{2}\sin^{3}(ak_0)},
\label{eq:dphi}\\[2pt]
R(k_0)&=|\tau(k_0)|\;e^{-\lambda^{2}/4\Delta_0^{2}},
\qquad \lambda=\frac{g/2t}{\sin^{2}(ak_0)} .
\label{eq:R}
\end{empheq}
The first term of Eq.~\eqref{eq:dphi} is the on-shell phase of
Eq.~\eqref{eq:Phi}. The second is the curvature correction
$\tfrac12\sigma_k^{2}\Phi''(k_0)$, and it is the first place the wavepacket
width enters the phase at all. Notwithstanding its smallness, of order one
percent for $\Delta_0=7$, it is the leading signature of the non-rigid
pointer and it is absent from any ideal von Neumann treatment.

Turning to Eq.~\eqref{eq:R}, the two factors correspond to two distinct
decoherence channels. The exponential is the pointer-overlap term familiar
from the continuum scheme, controlled by the Wigner displacement $\lambda$.
The prefactor $|\tau(k_0)|$ is elastic backscattering. For our parameters
the second dominates the first by between one and two orders of magnitude,
so that
\begin{equation}
1-R\simeq\tfrac12\,|r(k_0)|^{2}
=\frac{1}{2}\Big(\frac{g}{2t}\Big)^{2}\frac{1}{\sin^{2}(ak_0)} .
\label{eq:Rsimple}
\end{equation}
Coherence is therefore lost predominantly to the reflected channel and not
to pointer distinguishability. That identification explains why $R$ is
second order in $g$ while $\Delta\phi$ is first order, and it is consistent
with the entropy scaling reported in Sec.~III~B.

Collecting the results, the post-measurement Bloch vector is
\begin{equation}
\vec q=
\begin{pmatrix}
R\sin\theta_b\cos(\phi_0-\Delta\phi)\\
R\sin\theta_b\sin(\phi_0-\Delta\phi)\\
\cos\theta_b
\end{pmatrix}.
\label{eq:bloch}
\end{equation}
The kick retards the azimuth by $\Delta\phi$ and contracts the equatorial
projection by $R$, leaving the polar angle fixed. The sense of the rotation
follows the sign of $g$: a reduced bond amplitude, i.e.\ $g>0$, retards the
transmitted branch and the coherence lags, whereas an enhanced bond
advances it.

%---------------------------------------------------------------------
\subsection{Contrast with the ideal scheme, and validation}
\label{app:vn:contrast}

The ideal von Neumann coupling $\Phi(k)=-\lambda k$ has zero curvature.
Consequently its displacement is rigid and its phase is locked to the
carrier momentum, $\Delta\phi=\lambda k_0$, which grows with $k_0$. The
lattice QPC behaves oppositely: Eq.~\eqref{eq:dphi} decreases
monotonically with $k_0$ and vanishes at the band centre. Hence the
discrepancy is not a matter of coefficients. It reflects the fact that a
bond defect is a curved scatterer, whereas the ideal pointer coupling is by
construction flat. Our model is beyond von Neumann in this precise and
quantifiable sense: $\Phi''\neq0$.

We validate Eqs.~\eqref{eq:dphi} and \eqref{eq:R} against exact
time-dependent numerics on Eq.~(1) with $\Omega=0$, propagating the packet
through the bond and comparing the conditionally scattered state with the
freely evolved one. Table~\ref{tab:vn} collects the comparison for
$\Delta_0=7$ and $N=700$.

\begin{table}[t]
\caption{(a) Backaction phase and (b) coherence factor from exact
time-dependent numerics compared with Eqs.~\eqref{eq:dphi} and
\eqref{eq:R}, for $\Delta_0=7$, $N=700$, $\Omega=0$.}
\label{tab:vn}
\begin{ruledtabular}
\begin{tabular}{cc cc c cc c}
$g$ & $k_0/\pi$ & $\Delta\phi_{\rm ED}$ & Eq.~\eqref{eq:dphi} & ratio
& $1-R_{\rm ED}$ & Eq.~\eqref{eq:R} & ratio\\
\hline
0.100 & 0.30 & 0.037803 & 0.037802 & 1.0000 & $2.094\times10^{-3}$ & $2.035\times10^{-3}$ & 1.029\\
0.100 & 0.40 & 0.016841 & 0.016835 & 1.0003 & $1.489\times10^{-3}$ & $1.468\times10^{-3}$ & 1.015\\
0.050 & 0.30 & 0.018680 & 0.018674 & 1.0003 & $5.11\times10^{-4}$  & $4.97\times10^{-4}$  & 1.029\\
0.050 & 0.40 & 0.008319 & 0.008316 & 1.0004 & $3.63\times10^{-4}$  & $3.58\times10^{-4}$  & 1.014\\
0.025 & 0.30 & 0.009284 & 0.009280 & 1.0004 & $1.26\times10^{-4}$  & $1.23\times10^{-4}$  & 1.029\\
0.025 & 0.40 & 0.004134 & 0.004133 & 1.0003 & $9.0\times10^{-5}$   & $8.8\times10^{-5}$   & 1.014\\
\end{tabular}
\end{ruledtabular}
\end{table}

The phase agrees to better than one part in $10^{3}$ across the
perturbative range once the curvature term is retained. Dropping that term
degrades the agreement to the percent level, which confirms its
interpretation. The residual few-percent offset in $R$ is the next cumulant
and scales as $\sigma_k^{4}$, as expected.

%---------------------------------------------------------------------
\subsection{Regime of validity}
\label{app:vn:validity}

Three inequalities delimit the treatment. First, the kick picture on the
detector side requires no Brillouin-zone wrapping, $\lambda\ll\pi/k_0$.
Second, the perturbative expansion of Eq.~\eqref{eq:tauexact} requires
$g\ll2t$. Third, and separately, the collapse of the interaction picture on
the qubit side requires
\begin{equation}
\Omega\,\delta\tau\ll1 ,
\label{eq:slow}
\end{equation}
which follows from
$[\hat h_I,\hat H_0]=-i\Omega\,\hat\Phi(\hat k)\otimes\hat\sigma_y/\delta\tau$,
obtained using $\ket{1}\bra{1}=\tfrac12(\mathbb{1}-\hat\sigma_z)$ and
$[\hat\Phi(\hat k),\hat H_D]=0$. Equation~\eqref{eq:slow} is precisely the
slow-qubit condition of regime~[I], which identifies the slow-qubit limit
with the non-demolition limit. Outside it the kick picture fails on the
qubit side, the qubit rotates during the encounter, and the resonant
inelastic channel of Sec.~III~C opens. That breakdown is the origin of the
enhanced measurement band, and it is analysed perturbatively in
Appendix~C.
\end{document}